% ======================================================================
% 正文 第1章
% ======================================================================
\clearpage
\pagenumbering{arabic}
\setcounter{page}{1}

\pagestyle{fancy}
\fancyhf{}
\lhead{\zihao{-5}\songti \buildunit\project\reporttype}
\rhead{\zihao{-5}\songti \reportno}
\lfoot{\zihao{-5}\songti \company\quad 编制}
\rfoot{\zihao{-5}\songti 第\thepage\ 页 \quad 共\pageref{LastPage} 页}
\renewcommand{\headrulewidth}{0pt}
\renewcommand{\footrulewidth}{0pt}

\fancypagestyle{plain}{%
    \fancyhf{}%
    \lhead{\zihao{-5}\songti \buildunit\project\reporttype}%
    \rhead{\zihao{-5}\songti \reportno}%
    \lfoot{\zihao{-5}\songti \company\quad 编制}%
    \rfoot{\zihao{-5}\songti 第\thepage\ 页 \quad 共\pageref{LastPage} 页}%
    \renewcommand{\headrulewidth}{0pt}%
    \renewcommand{\footrulewidth}{0pt}%
}

\zihao{-4}\songti

\chapter{概述}

\section{项目背景}

萍乡市人民医院始建于 1928 年（原名普爱医院），1960 年更名，三甲综合医院，赣南医科大学非直属附属医院。2024 年挂号 121.57 万人次、出院 8.75 万人次、手术 3.9 万例，服务湘赣边约 1300 万人口；占地 218109㎡，建筑面积 172827㎡，开放床位 1728 张。医院持有《医疗机构执业许可证》《放射诊疗许可证》，许可含后装治疗等项目。

原机房山东新华 XHDR-18B 后装机 2017 年投用，设备老化损坏；医院拟淘汰旧设备，\textbf{新购置后装机（源}$^{192}\text{Ir}$\textbf{，活度}$3.7×10^{11}\text{Bq}(10\text{Ci})$\textbf{），原址机房使用}，定位沿用原有 SOMATOM go.Sim 大孔径 CT 模拟机。

依据《职业病防治法》《放射诊疗管理规定》，2025 年 6 月医院委托江西辐射剂量检测院开展本项目放射性职业病危害预评价。

\section{评价目的}

\begin{enumerate}
\item 落实法律法规要求，履行改建项目预评价法定程序；
\item 从选址、辐射危害、防护设施、监测、应急、管理制度论证项目放射防护可行性；
\item 为医院优化防护设计、行政部门审批提供技术依据。
\end{enumerate}

\section{评价范围}

\begin{table}[H]
\centering
\begin{tabular}{|c|c|c|}
\hline
范围分类 & 评价内容 & 备注 \\
\hline
辐射源范围 & 新购 1 台后装机（内置$^{192}\text{Ir}$密封源，近距离后装治疗） & 设备参数见 3.1 章节 \\
\hline
场所范围 & 设备使用区域、配套辅助用房 & 详见第 4 章 \\
\hline
防护安全设施 & 项目配套全部放射防护、安全设施 & / \\
\hline
人员范围 & 放射作业人员、可进入控制 / 监督区公众 & / \\
\hline
\end{tabular}
\end{table}


\begin{quote}
备注：不含土建施工阶段职业病危害评价；设备、防护重大变更需重新评价。
\end{quote}

\section{评价内容}

\begin{enumerate}
\item 项目辐射源项识别分析；
\item 放射性及伴生职业病危害因素识别；
\item 放射防护设施设计符合性评价；
\item 辐射监测方案核查；
\item 事故应急处置体系评价；
\item 放射防护管理、职业健康管理评价。
\end{enumerate}

\section{评价依据}

\subsection{法律法规、规章文件}

\subsubsection{法律}

\begin{enumerate}
\item 《职业病防治法》（2018 第四次修订，主席令 24 号）
\item 《放射性污染防治法》（2003.10.1，主席令 6 号）
\end{enumerate}

\subsubsection{行政法规}

\begin{enumerate}
\item 《放射性同位素与射线装置安全和防护条例》（国务院 449 号令，2019 修订）
\item 《医疗机构管理条例》（国务院 149 号令，2022 修订）
\item 《医疗器械监督管理条例》（国务院 276 号令，2020 修订）
\end{enumerate}

\subsubsection{部门规章}

\begin{enumerate}
\item 《放射诊疗管理规定》（卫生部 46 号令，2016 修改）
\item 《放射工作人员职业健康管理办法》（卫生部 55 号令）
\item 《医疗器械临床使用管理办法》（卫健委 8 号令）
\end{enumerate}

\subsubsection{规范性文件（国家 + 江西省）}

国家卫健委、原卫计委、江西省卫健委相关放射防护、培训、监管文件（卫监督发〔2012〕25 号、国卫办监督发〔2016〕38 号、赣卫监督发〔2019〕5/62 号等）。

\subsection{采用国家标准}

GB 18871-2002、GBZ 98-2020、GBZ 121-2020、GBZ 128-2019、GBZ 158-2003、GBZ/T 149-2015、GBZ/T 181-2024、GBZ/T201.1/201.3、GBZ/T220.2-2009、WS/T328-2011、WS262-2017。

\subsection{基础资料}

项目委托书（2025.6.10）、医院图纸、防护方案、人员资料、现有管理制度。

\subsection{参考资料}

ICRP103 报告、放射防护专业手册、国标指南、全国天然放射性水平调查报告。

\section{评价目标}

\subsection{辐射防护三原则：正当性、最优化、个人剂量限值}

\begin{table}[H]
\centering
\begin{tabular}{|c|c|}
\hline
受照类别 & 年有效剂量国标限值 \\
\hline
放射工作人员 & 连续 5 年年平均≤20mSv；单年度≤50mSv \\
\hline
公众 & 连续 5 年年平均≤1mSv；单年度≤5mSv \\
\hline
\end{tabular}
\end{table}


\subsection{医院内控管理目标}

\begin{table}[H]
\centering
\begin{tabular}{|c|c|}
\hline
人员分类 & 年有效剂量管控目标值 \\
\hline
放射工作人员 & 5mSv（国标 1/4） \\
\hline
周边公众 & 0.25mSv（国标 1/4） \\
\hline
\end{tabular}
\end{table}


\subsection{机房外剂量控制（GBZ121-2020）}

\begin{enumerate}
\item 居留因子 T＞1/2 区域：$\dot{H}_{c,max}\le2.5\mu\text{Sv/h}$；
\item 居留因子 T≤1/2 区域：$\dot{H}_{c,max}\le10\mu\text{Sv/h}$；
\item 实际控制值取理论计算值与上限值较小值。
\end{enumerate}

计算公式：$H_{c}\leqH_{e}/(t×U×T)$

\section{评价方法}

\begin{enumerate}
\item \textbf{检查表法}：对照法规标准逐项核查设计、设施、管理制度符合性；
\item \textbf{理论计算法}：屏蔽厚度、墙外辐射剂量理论核算，对标限值；
\item \textbf{现场调查法}：实地踏勘院区布局、机房现状、现有防护与管理。
\end{enumerate}

\section{评价程序}

准备阶段：受托→收集资料→现场初勘→确定评价单元→编制评价方案；
实施阶段：工程分析、源项识别、防护设施评价；
收尾阶段：编制初稿→内审→专家评审→修改→出具报批稿。
执行三稿（初稿、送审、报批）、三审制度。

\section{质量控制措施}

合同评审→组建项目组→资料收集审核→区分一般 / 严重类项目编制方案→现场勘查→编制报告→分级内审→专家评审→定稿归档。
